% ECONOMICS_PROBLEM: {"id": "L94-388", "title": "Electricity hedging and demand incentives", "jel_code": "L94", "source_id": "OP-0438"}
% !TeX program = xelatex
\documentclass[11pt,a4paper]{article}
\usepackage[margin=23mm]{geometry}
\usepackage{fontspec}
\setmainfont{Latin Modern Roman}
\usepackage{amsmath,amssymb,mathtools}
\usepackage{xcolor,enumitem,booktabs,longtable,array,xurl,hyperref}
\definecolor{muted}{HTML}{53636C}
\hypersetup{colorlinks=true,urlcolor=blue,linkcolor=blue}
\setlength{\parindent}{0pt}
\setlength{\parskip}{5pt}
\newcommand{\R}{\mathbb R}
\newcommand{\E}{\mathbb E}
\newcommand{\N}{\mathbb N}
\newcommand{\ind}{\mathbf 1}
\newcommand{\Var}{\operatorname{Var}}
\newcommand{\Cov}{\operatorname{Cov}}
\newcommand{\supp}{\operatorname{supp}}
\DeclareMathOperator*{\argmin}{arg\,min}
\DeclareMathOperator*{\argmax}{arg\,max}
\newcommand{\entrymeta}[3]{{\sffamily\small\color{muted}\textbf{\texttt{#1}}\quad|\quad #2\par #3\par}}
\newcommand{\Problemref}[1]{\texttt{#1}}
\newcommand{\JELref}[2]{\texttt{#2} (JEL \texttt{#1})}
\begin{document}
\section*{Electricity hedging and demand incentives}
\textbf{Problem ID:} \texttt{L94-388}\quad \textbf{JEL:} \texttt{L94}\par
% BEGIN ECONOMICS STATEMENT
\label{problem:OP-0438}
{\sffamily\small\color{muted}Original subject group: Climate, environment and energy\par}
\label{definitions:problem:30007720}
\textbf{Audit classification:} Published research agenda. Restricted otherwise vacuous free hedging and made spot-marginal retail incentives explicit.
\subsection*{Definitions and precise research target}
\textit{Editorial formalization.} Consider a wholesale electricity market with hourly spot prices \(p_t\), random demand, and uncertain renewable output. A retailer serves customers whose demand \(q_t(p_t,z)\) depends on the spot price and predetermined state \(z\). The retailer may buy a long-term hedge \(h\in\mathcal H\), where the allowed contract set \(\mathcal H\) specifies strike prices, quantities, settlement rules, and collateral requirements. Let \(\Pi(h)\) be retailer annual profit and \(C(h)\) total system resource cost after optimal demand and investment responses. Require finite second moments of profits and a finite lower bound on resource costs. For a stated variance tolerance \(\bar v\), determine whether there exists a feasible contract satisfying
\[\operatorname{Var}(\Pi(h))\leq\bar v\quad\text{and}\quad C(h)-\inf_{g\in\mathcal H}C(g)\leq\eta,\]
where \(\eta\geq0\) is a prespecified efficiency tolerance. Require customers to retain the marginal incentive to reduce consumption in scarce hours, and define how retail tariffs implement that requirement. Characterize the trade-off between hedging, credit risk, demand incentives, and investment incentives under the chosen market model. Compare mandatory procurement with voluntary contracts and specify enforcement and forecast errors.

Specify retail bill \(B_h(q,z)\), not merely wholesale prices. Preserving the scarcity incentive means \(\partial B_h/\partial q_t=p_t\) for each hour wherever the derivative exists, or equal corresponding marginal opportunity costs when taxes and network charges are included. A fixed-quantity financial hedge with payoff \(H_h(p)\) independent of customer marginal consumption can preserve that condition; a volume-indexed subsidy need not. Define \(\Pi(h,\omega)=\sum_t[B_h(q_t,z)-p_tq_t]+H_h(p)-\operatorname{premium}(h)-\operatorname{othercost}(h)\), plus defaults/collateral consistently. Set \(C(h)=E[C_{\rm resource}(h,\omega)]\) and assume \(\Pi(h)\in L^2\). If \(\mathcal H\) is unrestricted or contains free unlimited insurance, the proposed existence test is vacuous; impose a finite or compact, fairly priced implementable menu and a valid counterparty budget. All expectations use a stated baseline population and probability law. Identification means every admissible data-generating model with the same observed-data law yields the same displayed target; otherwise report the identified set. Exact equations, horizons and assumptions are editorial, rather than source-given canonical conjectures.
\subsection*{Variants and assumption changes}
\begin{enumerate}
\item \textbf{Editorial credit variant.} Add \(P(\Pi(h)<-c_0)\leq\alpha\) and collateral caps. A low-variance contract may fail this solvency condition. Fix the loss threshold \(c_0\geq0\) and failure probability \(\alpha\in(0,1)\) before comparison.
\item \textbf{Editorial tariff variant.} Compare \(B_h(q)=\sum_tp_tq_t+f_h\) with \(B_h(q)=\bar p\sum_tq_t+f_h\) under identical financial menus. A fixed fee cannot repair an incorrect marginal consumption incentive.
\end{enumerate}
\subsection*{References and source audit}
\textbf{1.} Explicit agenda on forward obligations and real-time demand incentives; source does not solve proposed constrained designs. Locator: Report 24-09, June 2024; Section 5, printed pp.38–40. Full text or primary publication page inspected. URL: \url{https://www.rff.org/documents/4511/RFF_Report_24-09.pdf}.
\begin{enumerate}\raggedright
\item\hypertarget{definitions:source:30007720:1}{} Chiara Lo Prete; Karen Palmer; Molly Robertson. 2024. \emph{Time for a Market Upgrade? A Review of Wholesale Electricity Market Designs for the Future}. Report 24-09, June 2024; Section 5, printed pp.38–40.
\url{https://www.rff.org/documents/4511/RFF_Report_24-09.pdf}.
\end{enumerate}

\subsection*{Common conventions: Source mathematical conventions}

These conventions apply to each entry unless explicitly replaced.

\textbf{Models and identification.} A primitive model \(m\) specifies all probability laws, feasible choices, information, equilibrium and measurement rules used by its target. The admissible class \(\mathcal M\) is fixed independently of outcomes. If \(P_O(m)\) is its observable probability law and \(\tau(m)\) the target, its identified set at observable law \(P\) is
\[
 \mathcal I_\tau(P)=\{\tau(m):m\in\mathcal M,\ P_O(m)=P\}.
\]
Point identification means that this set is a singleton. An empty set signals incompatibility of data and assumptions. Coordinate bounds need not describe the joint set. Bounds are sharp only if every retained target is attainable or arbitrarily approachable by compatible models. Finite bounds require explicit restrictions on unobserved quantities; unrestricted confounding, hidden wealth, extrapolation or arbitrary equilibrium selection can yield uninformative sets.

\textbf{Causal interventions.} A potential outcome \(Y_i(p)\) is the outcome under a complete policy regime \(p\), including timing, financing, information and market spillovers. The target population and units are fixed before treatment. With interference, use the complete assignment vector or a justified exposure mapping; individual no-interference notation alone cannot identify an economy-wide intervention. A difference of mean potential outcomes does not require the cross-world joint distribution. Conditioning on post-treatment migration, survival, enrollment or employment changes the estimand and needs separate justification.

\textbf{Statistical statements.} An estimator, confidence set, forecast or test specifies a sampling law, model class and loss/coverage criterion. Uniform \((1-\alpha)\)-coverage means \(\inf_{m\in\mathcal M}\Pr_m\{\tau(m)\in C_n\}\geq1-\alpha\), or an explicitly asymptotic version. The whole parameter space trivially has coverage, so informativeness requires a stated width, risk or power objective. A forecasting improvement is relative to a named benchmark, information set and evaluation population.

\textbf{Equilibrium and optimization.} An equilibrium comprises feasible plans, optimality, beliefs where needed, budget constraints and all market/resource clearing. The primitives or a stated benchmark define each correspondence \(\mathcal E(p;m)\). Nonemptiness is assumed or proved, never inferred just from compact policy sets. A selected-equilibrium target uses a declared selection rule; a robust target quantifies over all permitted equilibria. Use suprema or infima unless attainment is established. An \(\varepsilon\)-optimal policy is feasible and within \(\varepsilon\) of the finite optimum value. Report an empty feasible set separately.

\textbf{Welfare and accounting.} Normative utility, interpersonal normalization, social weights, numeraire, discounting and the use of public receipts are primitives. Behavioral utility need not coincide with the welfare criterion. Household welfare includes ownership income and rents once; adding profits again to compensating variation generally double-counts them. A monetary equivalent is defined by an explicit transfer and price convention, with monotonicity and range conditions. Average effects, mechanism contributions, transfers and welfare losses are different targets. Infinite sums require convergence and dynamic feasible plans require terminal or no-Ponzi conditions.

\textbf{Source versions.} A working-paper date, revision date and journal publication year are distinguished. A DOI for a journal version is not automatically the DOI of an earlier manuscript. Locators refer to the stated version and printed pages where specified. A metadata-only check does not verify a claim about a full-text conclusion. Every entry's source subsection narrows the provenance claim accordingly.
% END ECONOMICS STATEMENT
\end{document}
